% RESS Sections 1-2 (R1 approved 30-09-2026; Section 1 not locked until the TT 11 reject reasons are confirmed).
\section{Introduction}\label{sec:intro}

A reliability demonstration test (RDT) accepts a product when it passes a planned test; for attribute (pass/fail) data
the classical plan is the zero-failure, or success-run, test~\cite{hanley1983,meeker2017}. If each demand fails
independently with probability $p$, a unit with $p>\varepsilon$ passes $n$ demands with probability at most
$(1-\varepsilon)^n$, so passing $n_0(\varepsilon,\delta)=\lceil\ln\delta/\ln(1-\varepsilon)\rceil$ demands without
failure demonstrates $p\le\varepsilon$ at confidence $1-\delta$: \nNzero{} demands at $\varepsilon=\delta=5\%$. Zero-failure
planning remains common industrial practice~\cite{yoon2025bayesian}, and its sample-size burden is a recurring motive for
test-plan research in this journal, from success--failure verification tests~\cite{ding2024truncated} to qualification
plans that borrow auxiliary information~\cite{wang2024monoprop,tan2025weibull}.

We consider a learned perception component, a fixed object detector on a fixed camera that samples frames
periodically, as the unit under test. A demand is the passage of an object (an \emph{event}); a failure is an event on
which no sampled frame yields a detection. Such components are increasingly assessed with reliability-engineering
tools~\cite{wen2025nversion,paterson2025amlas,aghazadeh2026hipllm}. For them, the demands are rare and every new
installation is a new unit: \nNzero{} failure-free events must be collected per camera, and for road vehicles the
exposure needed to demonstrate rare-failure reliability by field operation alone is prohibitive~\cite{kalra2016}.

Simulation is the obvious remedy. Digital twins and simulators already supply reliability evidence in engineering
practice~\cite{tao2024subsea,ye2023structuraldt,wang2026fusion}, and virtual scenario testing is a leading proposed alternative to mileage-based demonstration of automated vehicles~\cite{riedmaier2020,zhao2025statfoundation}. For perception,
however, the simulation-to-reality gap is measurable and has been measured under different environmental conditions~\cite{reway2020simgap,dieter2023drone}, and
performance in simulation need not predict performance in the field~\cite{kadian2020}. The question for a demonstration
test is therefore quantitative: \emph{how many of the $n_0$ real demands can simulated evidence replace, and what form of
guarantee survives when the evidence is transferred to an untested unit?}

This paper answers the question for zero-failure demonstration of a perception component.
\begin{enumerate}
\item \emph{No free credit.} Without an assumption that links the simulator to reality, no valid procedure can use
independent simulated data to shorten a zero-failure demonstration, for any simulator and any amount of simulation.
This makes precise, for success-run tests, a principle known from clinical
borrowing~\cite{koppschneider2020}, and relates to the limits on combining evidence in software
dependability~\cite{littlewood1993}.
\item \emph{The exact price of relaxing confidence.} Accepting error probability $\delta'>\delta$ when the simulator is
wrong buys at most, and exactly, $n_0(\varepsilon,\delta)-n_0(\varepsilon,\delta')$ demands: \ncExactA, \ncExactB,
\ncExactC{} and \ncExactD{} of the \nNzero{} for $\delta'=7.5$, 10, 15 and 20\%.
\item \emph{A paired-twin bound and a fleet-level transfer bound.} A synthetic twin of each real event bounds the real
failure probability by the twin failure probability plus a discordance rate, with no assumption; on one camera this
costs more real events than the plain test (\nCthreeA--\nCthreeC). Calibrated once across scenes, the discordance
transfers to an untested camera as a population-average bound whose price is the number of calibration scenes, not the
number of events.
\item \emph{An exchange rate on public benchmarks, and a named failure mode.} On \nMainScenes{} fixed-camera scenes the twin
as built saves no real event; a twin run to convergence would save at most \nSavedOP{} real events per camera, if the discordance calibrated with
paired twins also held for the twins of new cameras (not verified), at the
operating point fixed in advance, where real failures are rare (\nPmainOP) and the demonstration is therefore untested.
The transfer certificate stands at scene level but is withdrawn when the unit is the site (\nSiteSm{} sites).
When out-of-domain cameras are judged alone, they pass through: \nNightFC{} of \nNightScenes{} night cameras and, at a
fixed sampling phase, \nOODSJitfa{} of \nJitterScenes{} camera-jitter cameras are falsely accepted; we call this
\emph{out-of-domain pass-through}, with two named modes, night and camera jitter.
\end{enumerate}

\emph{Reliability demonstration tests in the zero-failure regime have no established rule for how much simulated
evidence may substitute for real evidence, nor a guarantee form that survives transfer to an untested unit; this paper
gives a finite-sample answer (none without an assumption; an exact price for relaxing confidence; a fleet-level
bound whose price is scene count) and measures it on \nMainScenes{} fixed-camera scenes.}

% =====================================================================================================
\section{Related work and research gap}\label{sec:related}

Table~\ref{tab:gap} positions the paper against four bodies of work.

% Table 1 (RESS): research gap. Row (e) numbers via numbers.tex.
\begin{table*}[t]
\centering
\caption{Research gap. Unit: what the guarantee covers. ZF: is the zero-failure (success-run) regime treated? Price: is the
real-evidence value of simulation quantified? Transfer: does the guarantee reach an untested unit? Link: assumption needed
between simulation and reality. Evidence: empirical basis.}
\label{tab:gap}
\scriptsize
\setlength{\tabcolsep}{2pt}
\begin{tabular}{@{}L{0.18\textwidth}L{0.10\textwidth}L{0.05\textwidth}L{0.16\textwidth}L{0.15\textwidth}L{0.15\textwidth}L{0.13\textwidth}@{}}
\toprule
Body of work & Unit & ZF & Price of simulation & Transfer & Sim--real link & Evidence \\
\midrule
(a) Reliability demonstration, zero-failure planning~\cite{hanley1983,baze1979,wilson2021assurance,zheng2023gamma,ding2024truncated,wang2024monoprop,zhao2018perfection}
 & product / unit & yes & through a prior or auxiliary data; not as real events saved & to the tested product; prior-system evidence handled conservatively & prior or similarity assumption (Bayesian); none (frequentist, no simulation) & case studies \\
(b) Simulation-assisted testing, digital-twin reliability~\cite{tao2024subsea,ye2023structuraldt,wang2026fusion,li2026tugs,riedmaier2021vvuq}
 & system & no & fidelity weighting in fusion & not addressed (extrapolation named as open) & fidelity or calibration model & engineering case studies \\
(c) Simulation-augmented evaluation of perception~\cite{reway2020simgap,dieter2023drone,amini2024translators,zhao2025statfoundation,angelopoulos2023ppi,fisch2024stratppi}
 & mean metric & no & gap measured; variance reduction for means & no & paired labels (rectifier) or none & perception benchmarks \\
(d) Finite-sample guarantees and assurance for learned perception~\cite{angelopoulos2025ltt,degrancey2022conformaldet,mei2025pwc,yuan2026conformal,paterson2025amlas}
 & per deployment & no & not addressed (real calibration data) & exchangeability of calibration and deployment & exchangeability & detection data sets \\
(e) This paper & fleet (population average); per-unit as a data requirement & yes & exact: none without a link; $n_0(\delta)-n_0(\delta')$ when relaxing; \nSavedOP{} events per camera at most & yes, priced by scene count ($\UCP(J,m)=\nMainUpop$, $J=\nMainJ$, $m=\nMainScenes$) & none for the twin bound; scenes exchangeable (i.i.d.) for transfer, as in (d) at scene level & \nMainScenes{} fixed-camera scenes, \nMainN{} events; failure mode measured \\
\bottomrule
\end{tabular}
\end{table*}


\emph{(a) Reliability demonstration and zero-failure test planning.} The zero-numerator bound and the rule of three are
classical~\cite{hanley1983,meeker2017}, and the Bayesian zero-failure (BAZE) test sets the required test time through a prior on the failure rate~\cite{baze1979}. Recent work designs sample sizes by assurance~\cite{wilson2021assurance}, studies the
limits of Bayesian zero-failure plans~\cite{yoon2025bayesian}, and publishes in this journal demonstration and acceptance
plans for degrading products~\cite{zheng2023gamma,zheng2023acceptance}, success--failure verification
tests~\cite{ding2024truncated} and qualification plans that use subsystem, Monte-Carlo or expert
information~\cite{wang2024monoprop,tan2025weibull}. For software and autonomous systems, conservative Bayesian inference uses prior knowledge, including evidence from previous similar systems, without optimistic bias~\cite{zhao2018perfection,zhao2019cbi,zhao2020operational},
and reusing the reliability of an older system as a prior is optimistic rather than
conservative~\cite{littlewood2020replacement}. Bayesian borrowing quantifies how much a prior is worth through its effective sample size~\cite{morita2008ess} and
discounts historical data through power or commensurate priors~\cite{ibrahim2000power,hobbs2011commensurate}, and
conservative frequentist bounds have been derived when the test profile differs from the operational
profile~\cite{bishop2017profiles}. \emph{What is missing:} a frequentist statement of how much the
zero-failure sample size can be reduced by simulated rather than prior-system evidence, and at what price.

\emph{(b) Simulation-assisted testing and digital-twin reliability.} Digital twins feed reliability models with simulated
operating data~\cite{tao2024subsea}, are validated against physical experiments~\cite{ye2023structuraldt}, and are
discussed as maintenance tools whose suitability is still open~\cite{dwight2026dtmaint}. Multi-fidelity fusion combines
scarce physical tests with abundant simulation for mission-reliability prediction, weighting the simulation by its
fidelity gap~\cite{wang2026fusion}, and simulators generate test scenarios for autonomous
systems~\cite{li2026tugs,riedmaier2020}. Model verification, validation and uncertainty quantification identify extrapolation to untested conditions as a major challenge~\cite{riedmaier2021vvuq}. \emph{What is missing:} a rule that
states which assumption about the simulator is needed before its output may reduce a zero-failure test, and a measured
sim--real discordance that makes the assumption checkable.

\emph{(c) Simulation-augmented evaluation of perception.} The simulation-to-reality gap of camera object detectors has been
measured across day, night, fog and rain~\cite{reway2020simgap} and for drone detection~\cite{dieter2023drone}; synthetic
images are used to test driving perception, with a distribution mismatch that biases test
results~\cite{amini2024translators}, and perception digital twins reproduce field data only
approximately~\cite{saad2025dtperception}. Statistical model checking evaluates perception in simulation
only~\cite{barbier2019smc}, simulation performance can fail to predict field performance~\cite{kadian2020}, and a
statistical foundation for aligning synthetic and real testing outcomes has been called for~\cite{zhao2025statfoundation}.
Closest to our question, Sim2Val uses simulator outputs as control variates to reduce the number of real samples needed
for a confidence bound on mean performance (CoRL 2025, arXiv preprint)~\cite{luo2025sim2val}, SCAPE combines limited real rollouts with
simulation proxies for policy evaluation (preprint)~\cite{zhu2026scape}, and multi-fidelity fusion weights simulation by
its fidelity gap in reliability prediction~\cite{wang2026fusion}; none treats the zero-failure regime.
Prediction-powered inference and its stratified variant combine many proxy outcomes with few real labels through a paired
rectifier~\cite{angelopoulos2023ppi,fisch2024stratppi}. \emph{What is missing:} these methods estimate a mean performance
or reduce its variance; none addresses the zero-failure regime, where the quantity to bound is a rare failure probability
and the rectifier itself would cost almost $n_0$ real events per unit.

\emph{(d) Finite-sample guarantees and assurance for learned perception.} Conformal prediction gives distribution-free finite-sample
guarantees~\cite{angelopoulos2023gentle}, calibrated risk control extends them to general losses~\cite{angelopoulos2025ltt},
and they have been applied to object detection~\cite{degrancey2022conformaldet}, keypoint-based pose
estimation~\cite{yang2023purse}, perception for safe navigation~\cite{mei2025pwc} and, in this journal, anomaly detection in
cyber-physical systems~\cite{yuan2026conformal}. Reliability of machine-learning components is also modelled directly,
through multi-version architectures~\cite{wen2025nversion}, safety-assurance cases for object
detectors~\cite{paterson2025amlas} and failure-free-operation definitions under an operational
profile~\cite{aghazadeh2026hipllm}. \emph{What is missing:} these guarantees are calibrated on real data from the
deployment distribution; they do not say how much of that calibration data may be synthetic, nor how a guarantee
transfers to a unit whose own data are absent.
